
\chapter{جمع‌بندی و کارهای آتی}
\label{chap:conclusions}

\section{جمع‌بندی}
\label{sec:c-path}

این پژوهش به بررسی کاهش شکاف دقت بخش‌بندی چندعضوی شکم میان پویش بدون‌تزریق و پویش دارای ماده‌ی حاجب، و ارائه‌ی یک زنجیره‌ی کارآمد برای این منظور پرداخته است. برای دستیابی به این هدف، ساختار پایان‌نامه در شش فصل تنظیم شده است که هر یک به بخش خاصی از موضوع اختصاص دارد.

در فصل اول، ابتدا اهمیت بالینی مسئله معرفی شده است. پویش بدون‌تزریق برای گروهی از بیماران، از جمله بیماران ترومایی و آن‌ها که به دلیل نارسایی کلیوی امکان تزریق ماده‌ی حاجب ندارند، تنها تصویر در دسترس است و با این حال دقت بخش‌بندی خودکار روی آن به‌طور نظام‌مند پایین‌تر است. سپس مسئله به چهار زیرمسئله‌ی مشخص تجزیه شده و چالش‌ها و اهداف پژوهش بیان شده است.

در فصل دوم، پیشینه‌ی پژوهشی در سه حوزه‌ی شکاف بخش‌بندی، هم‌ترازسازی بین‌پویشی و هم‌نهادی پویش دارای ماده‌ی حاجب مرور شده است. این فصل نشان داد که شکاف بخش‌بندی وجود دارد اما یکنواخت نیست و در ساختارهای کوچک و کم‌مرز مانند پانکراس بیشتر است. همچنین معیارهای هم‌ترازسازی مبتنی‌بر شدت روشنایی در برابر تزریق ماده‌ی حاجب شکننده‌اند و همه‌ی روش‌های موجود معیار تراز را روی کل میدان دید محاسبه می‌کنند. در حوزه‌ی هم‌نهادی نیز تقریباً هیچ روشی خروجی خود را بر پویش و سطح تقویت شرطی نمی‌کند و کیفیت هم‌ترازی جفت‌های آموزشی گزارش نشده است.

در فصل سوم، مرور نظام‌مند مجموعه‌داده‌های عمومی سی‌تی شکمی گزارش شده است. این مرور نشان داد که \عدد{۴۹٫۵} درصد حجم ظاهری داده‌ی عمومی بازاستفاده‌ی موارد پیشین است، \عدد{۷۵٫۳} درصد مجموعه‌داده‌ها از آمریکای شمالی و اروپا می‌آیند و \عدد{۱۲} مورد از \عدد{۲۲} مجموعه‌داده‌ی بزرگ در بالاترین رده‌ی مخاطره‌ی سوگیری قرار دارند. سهم پویش بدون‌تزریق از داده‌ی برچسب‌خورده نیز در بهترین حالت حدود \عدد{۲۲} درصد است؛ رقمی که پیش از این در هیچ منبعی گزارش نشده بود و مسیر آموزش مستقیم روی این پویش را عملاً می‌بندد. از همین رو جفت‌های آموزشی باید از داده‌ی چندپویشی ساخته شوند.

در فصل چهارم، روش پیشنهادی ارائه شده است که از سه مؤلفه‌ی متوالی تشکیل می‌شود. نخست یک دسته‌بند پویش که از میان روش‌های منتشرشده برگزیده شده و پویش هر سری را از میانه‌ی واحد هانسفیلد اندام‌ها تشخیص می‌دهد. دوم یک زنجیره‌ی هم‌ترازسازی صلب چندگذره که به‌جای وزن‌دهی، تبدیل را تنها از ساختارهایی برآورد می‌کند که میان دو پویش جابه‌جا نمی‌شوند و بخش‌بندی‌شان پایدار است. سوم شبکه‌ی \lr{PLC-Net} که پویش و سطح تقویت هدف را به‌عنوان ورودی می‌پذیرد و با یک مدل واحد هر دو پویش را تولید می‌کند.

در فصل پنجم، نتایج آزمایش‌ها روی مجموعه‌داده‌ی \lr{VinDr-Multiphase} گزارش شده است. دسته‌بند پویش روی بیماران ناشناخته به دقت \عدد{۰٫۹۸۴} رسید. در هم‌ترازسازی، پیش‌پردازش پیشنهادی شکست‌ها و خطاهای بزرگ هر چهار روش هم‌ترازسازی موجود را به‌طور معنادار کاهش داد و، روی نقاب اندام‌ها، تابع آناتومیک ساده از تابع لبه بهتر بود. \lr{PLC-Net} نیز دایس بخش‌بندی ورید باب و طحالی را در هر \عدد{۲۰} مورد آزمون به‌اندازه‌ی \عدد{۰٫۱۶۶} افزایش داد. سودمندی پایین‌دستی آن با \lr{ResViT} هم‌ارز است (در حاشیه‌ی \عدد{۰٫۰۲} دایس)، اما مولد آن حدود \عدد{۹} برابر پراسنجه‌ی کمتر و حدود \عدد{۷} برابر حافظه‌ی کمتر دارد؛ شمار عمل به‌ازای هر برش در ارزیابی دو مدل برابر است، چون \lr{PLC-Net} هر برش را در وصله‌های هم‌پوشان می‌پردازد، و زمان واقعی اندازه‌گیری نشد. با این حال اثر آن بر ترکیب از پیش تعیین‌شده‌ی اندام‌های وابسته به ماده‌ی حاجب، به دلیل افت متمرکز کیسه‌ی صفرا، معنادار نبود. بنابراین شکاف بخش‌بندی یکنواخت نیست و مداخله نیز یکنواخت اثر نمی‌گذارد: در ساختارهایی که بدون ماده‌ی حاجب مرز قابل‌تشخیص دارند، مانند کبد، طحال و کلیه‌ها، دقت بخش‌بندی از پیش بالاست و جایی برای بهبود نمی‌ماند، و سود قابل‌دفاع در ساختاری متمرکز است که بدون ماده‌ی حاجب از پارانشیم اطراف جدا نمی‌شود. افزون بر آن، ارزیابی معیارها نشان داد که در میان معیارهای تصویری، تنها خطای واحد هانسفیلد در سطح اندام با سودمندی پایین‌دستی هم‌بسته است و معیارهای متعارف \lr{PSNR}، \lr{SSIM}، \lr{FID} و \lr{LPIPS} چنین رابطه‌ای نشان نمی‌دهند.

لازم به ذکر است که همه‌ی آزمایش‌ها روی یک مجموعه‌داده‌ی عمومی و تک‌مرکزی، با مجموعه‌های آزمون کوچک و یک بذر تصادفی انجام شده‌اند؛ دامنه‌ی اعتبار نتایج در بخش~\ref{sec:e-summary} آمده است و بیشتر کارهای آتی زیر همین محدودیت‌ها را هدف می‌گیرند. در پایان، نتایج مرور نظام‌مند این پژوهش به‌طور مستقل نیز منتشر شده است~\cite{danaei2026abdomenctreview}.

\section{کارهای آتی}
\label{sec:c-future}

در مسیر پیشرفت در زمینه‌ی کاهش شکاف بخش‌بندی در پویش بدون‌تزریق، \lr{PLC-Net} و زنجیره‌ی هم‌ترازسازی پیشنهادی همچنان کمبودهایی دارند و نقاط بهبود متعددی وجود دارد که می‌توان با پرداختن به آن‌ها به نتایج دقیق‌تر و کاربردی‌تری دست یافت. در ادامه، مهم‌ترین مسیرهای پژوهشی برای آینده مورد بررسی قرار می‌گیرند:

\subsection*{بهبود عملکرد در کیسه‌ی صفرا}

این تنها اندامی است که همه‌ی بازوهای هم‌نهاده در آن از بازوی بدون‌تزریق بدتر عمل می‌کنند، و همین افت است که اثر مثبت ورید باب را در نقطه‌ی پایانی تأییدی خنثی می‌کند. شکست از جنس تشخیص نیست بلکه کم‌برآوردی نظام‌مند حجم است. نخستین آزمون مستقیم، مقایسه‌ی \lr{HU} درون کیسه‌ی صفرا در تصویر هم‌نهاده و واقعی است که در این پایان‌نامه انجام نشد. اگر فرضیه‌ی روشن‌شدن حفره تأیید شود، خارج‌کردن کیسه‌ی صفرا از بردار سطح یا وزن‌دهی جداگانه‌ی آن در تابع زیان، گام بعدی است.

\subsection*{جداسازی دو علامت شرط}

متعارف‌سازی بردار سطح باید \emph{درون هر پویش} انجام شود و مدل با همین بردار بازآموزی شود؛ یک طرح صریح حذف تصادفی شرط در آموزش نیز شبکه را وادار می‌کند به هر دو علامت مستقلاً تکیه کند. این تغییر، هم دستگیره‌ی سطح شریانی را که اکنون کار نمی‌کند فعال می‌کند و هم کنترل‌پذیری را از خوانش مبهم فعلی بیرون می‌آورد.

\subsection*{ارزیابی برون‌مرکزی}

همه‌ی ادعاهای این پایان‌نامه روی یک مجموعه‌داده آزموده شده‌اند. آزمودن هر سه مؤلفه روی داده‌ی مرکز دوم، به‌ویژه با حضور پویش تأخیری، نخستین چیزی است که دامنه‌ی این ادعاها را روشن می‌کند. فصل~\ref{chap:survey} نشان داد کدام منابع عمومی برای این کار مناسب‌ترند.

\subsection*{آزمون مستقیم سود محدودکردن دامنه}

در آزمایش مقایسه‌ی توابع هدف، هر چهار تابع روی نقاب اندام‌ها حساب شدند و هیچ‌کدام روی کل تصویر اجرا نشد؛ پس سود محدودکردن دامنه آزموده نشده است. اجرای یک معیار یکسان (اطلاعات متقابل، همبستگی لبه‌ها یا \lr{MIND}) یک‌بار روی کل تصویر و یک‌بار روی نقاب، با همه‌چیز دیگر ثابت، این ادعا را از یک انتخاب طراحی به یک مقایسه‌ی تک‌متغیره تبدیل می‌کند.

\subsection*{تکمیل ارزیابی احتمالاتی}

نمونه‌برداری چندگانه از مدل انتشار و محاسبه‌ی امتیاز احتمال رتبه‌ای پیوسته و پوشش بازه، پیش‌بینی ثبت‌شده در بخش~\ref{sec:m-uncertainty} را آزمون‌پذیر می‌کند: این‌که یک نمونه‌ی صادق از توزیع شرطی، زیر یک معیار نقطه‌ای، باید حدود $\sqrt{2}$ برابر بدتر از میانگین شرطی امتیاز بگیرد.

\subsection*{مرجع انسانی برای ارزیابی پایین‌دستی}

حاشیه‌نویسی دستی حتی زیرمجموعه‌ی کوچکی از موارد آزمون، به‌ویژه روی ورید باب و کیسه‌ی صفرا، تنها راه بیرون‌آمدن از حلقه‌ای است که در آن تصویر هم‌نهاده با شبیه‌شدن به خروجی یک ابزار خودکار امتیاز می‌گیرد.

\subsection*{انتشار کد و جفت‌های هم‌ترازشده}

بخش~\ref{sec:synthesis} مستند کرد که در هسته‌ی مسئله‌ی هم‌نهادی چندپویشی شکمی تقریباً هیچ کد بازی وجود ندارد. انتشار زنجیره‌ی هم‌ترازسازی، که هیچ وزن آموخته‌شده‌ای ندارد و روی یک مجموعه‌داده‌ی عمومی اجرا شده است، به‌همراه شاخص‌های کیفیت هم‌ترازی هر جفت، کم‌هزینه‌ترین سهمی است که این پروژه می‌تواند به بازتولیدپذیری این حوزه بدهد.
